% ============================================================
%  2.1 – Thiết kế giao diện người dùng (UI/UX)  (1.5 điểm)
%  Trạng thái: ĐÃ HOÀN THIỆN
% ============================================================

\subsection{Phân tích người dùng}

UniJob phục vụ mô hình \textbf{C2C} (Campus-to-Campus) trong hệ sinh thái HCMUT, kết nối
hai nhóm sinh viên có nhu cầu bổ sung lẫn nhau. Qua khảo sát nhóm và tham khảo nhu cầu
thực tế từ cộng đồng sinh viên Bách Khoa, nhóm xác định \textbf{hai personas chính}:

\begin{table}[H]
\centering
\caption{Personas người dùng UniJob}
\renewcommand{\arraystretch}{1.2}
\begin{tabularx}{\textwidth}{|p{3.2cm}|X|X|}
\hline
\textbf{Tiêu chí} & \textbf{Persona A -- Người đăng task} & \textbf{Persona B -- Người tìm task} \\
\hline
\textbf{Đặc điểm} & SV năm 3--4, bận rộn với đồ án, cần outsource công việc nhỏ trong campus &
  SV năm 1--3, có kỹ năng chuyên biệt (lập trình, thiết kế, gia sư...), muốn kiếm thêm thu nhập \\
\hline
\textbf{Nhu cầu} & Tìm được người có kỹ năng phù hợp nhanh trong cộng đồng đáng tin cậy &
  Task phù hợp kỹ năng, địa điểm gần, thù lao rõ ràng, có thể tích lũy portfolio \\
\hline
\textbf{Điểm đau} & Không biết tìm ai trong mạng lưới; lo ngại về độ tin cậy và chất lượng &
  Thiếu kênh chính thống để offer kỹ năng; không có reputation để chứng minh năng lực \\
\hline
\textbf{Mục tiêu UX} & Đăng task nhanh ($<$2 phút), theo dõi ứng viên, xác nhận hoàn thành dễ dàng &
  Tìm và ứng tuyển task $<$3 bước, xây dựng hồ sơ uy tín, xuất CV Passport \\
\hline
\end{tabularx}
\end{table}

Hai personas này dẫn đến yêu cầu thiết kế cốt lõi: giao diện \textbf{dual-role} (cùng
người dùng vừa đăng vừa tìm task), hệ thống \textbf{uy tín hai chiều} (rating sau khi
hoàn thành), và tính năng \textbf{CV Passport} giúp sinh viên chuyển kinh nghiệm thực
tế thành chứng nhận có thể chia sẻ.

\subsection{Sitemap và luồng điều hướng}

Dựa trên phân tích personas, nhóm thiết kế sitemap tối giản gồm \textbf{6 màn hình chính},
tất cả truy cập từ Header navigation:

\begin{figure}[H]
\centering
\begin{tikzpicture}[
  box/.style={draw, rounded corners, minimum width=2.5cm, minimum height=0.8cm,
              fill=green!10, font=\small, align=center},
  sbox/.style={draw, rounded corners, minimum width=2.5cm, minimum height=0.8cm,
               fill=gray!10, font=\small, align=center},
  arr/.style={->, thick, >=stealth}
]
  \node[box, fill=green!30, minimum width=3cm] (home) at (0,0) {Trang chủ (Home)};

  \node[box]  (list)   at (-5.5,-2) {Danh sách\\việc làm};
  \node[box]  (create) at (-2,-2)   {Đăng\\việc làm};
  \node[box]  (prof)   at (2,-2)    {Hồ sơ\\cá nhân};
  \node[sbox] (login)  at (5.5,-2)  {Đăng nhập\\(Google OAuth)};

  \node[box]  (detail) at (-5.5,-4) {Chi tiết\\việc làm};
  \node[box]  (cv)     at (2,-4)    {Xuất CV\\Passport};

  \draw[arr] (home) -- (list);
  \draw[arr] (home) -- (create);
  \draw[arr] (home) -- (prof);
  \draw[arr] (home) -- (login);
  \draw[arr] (list) -- (detail);
  \draw[arr] (prof) -- (cv);

  \node[font=\scriptsize, gray] at (-1,-3) {\textit{(yêu cầu đăng nhập)}};
\end{tikzpicture}
\caption{Sitemap UniJob -- cấu trúc điều hướng 6 màn hình chính}
\end{figure}

Luồng người dùng chính theo persona:
\begin{itemize}
  \item \textbf{Tìm task:} Trang chủ \textrightarrow{} Danh sách \textrightarrow{}
    Chi tiết \textrightarrow{} Ứng tuyển (inline form) \textrightarrow{} Hồ sơ \textrightarrow{}
    Xuất CV Passport.
  \item \textbf{Đăng task:} Trang chủ \textrightarrow{} Đăng việc làm \textrightarrow{}
    Hồ sơ cá nhân (xem danh sách ứng viên, xác nhận hoàn thành).
\end{itemize}

\subsection{Tổng quan thiết kế}

UniJob được thiết kế với triết lý \textbf{mobile-first, minimalist}, ưu tiên luồng người
dùng rõ ràng và giảm tối đa số bước để hoàn thành một nhiệm vụ (đăng task hoặc ứng tuyển).
Toàn bộ giao diện được thiết kế trên Figma và triển khai bằng
\textbf{React 19 + Tailwind CSS v4}, đảm bảo tính nhất quán giữa bản thiết kế và sản phẩm
thực tế.

\subsection{Hệ thống màu sắc và Design Tokens}

Nền tảng sử dụng một hệ màu xanh lá (\textit{emerald}) làm màu chủ đạo, thể hiện sự tin cậy
và gần gũi với môi trường học thuật.

\begin{table}[H]
\centering
\caption{Design tokens màu sắc chính của UniJob}
\renewcommand{\arraystretch}{1.2}
\begin{tabularx}{\textwidth}{|p{5.2cm}|p{3.3cm}|X|}
\hline
\textbf{Token} & \textbf{Giá trị} & \textbf{Sử dụng} \\
\hline
\texttt{--color-primary}     & \texttt{\#10b981} (emerald-500) & CTA buttons, badges, highlights \\
\hline
\texttt{--color-primary-dk}  & \texttt{\#059669} (emerald-600) & Hover state, accent \\
\hline
\texttt{--color-background}  & \texttt{\#ffffff}               & Nền trang chính \\
\hline
\texttt{--color-border}      & \texttt{\#e5e7eb} (gray-200)    & Viền card, divider \\
\hline
\texttt{--color-muted}       & \texttt{\#f9fafb} (gray-50)     & Card background \\
\hline
\texttt{--color-foreground}  & \texttt{\#111827} (gray-900)    & Text chính \\
\hline
\texttt{--color-muted-foreground} & \texttt{\#6b7280} (gray-500) & Text phụ, placeholder \\
\hline
\end{tabularx}
\end{table}

\subsection{Cấu trúc màn hình và luồng người dùng}

\subsubsection{Màn hình trang chủ (Home)}

Trang chủ chia thành các section theo thứ tự ưu tiên thông tin:

\begin{enumerate}
  \item \textbf{Hero Section:} Banner gradient xanh lá với headline ``Kết nối công việc ngay
    trong campus'', nút CTA ``Tìm việc ngay'' và ``Bắt đầu miễn phí''. Sử dụng
    \texttt{framer-motion} cho animation fade-in khi tải trang.

  \item \textbf{Stats Bar:} 4 số liệu nổi bật (số task đang mở, số người dùng, tỷ lệ hoàn
    thành, đánh giá trung bình) hiển thị ngay dưới hero để xây dựng tin cậy ngay lập tức.

  \item \textbf{Đề xuất cho bạn (SuggestedJobs):} Chỉ hiển thị khi đăng nhập. Hiển thị
    lưới 3 cột gồm 6 công việc gợi ý dựa trên thuật toán AI Matching (xem mục 5.1).

  \item \textbf{How it works:} 3 bước trực quan (Đăng ký → Tìm task → Hoàn thành \&
    Đánh giá) với icon và mô tả ngắn.

  \item \textbf{CTA cuối trang:} Kêu gọi đăng ký hoặc đăng task lần cuối.
\end{enumerate}

\subsubsection{Màn hình danh sách công việc (JobList)}

\begin{itemize}
  \item Layout 2 cột: sidebar bộ lọc bên trái (categories, faculty, salary range, work mode,
    urgent flag) + danh sách card bên phải.
  \item Thanh tìm kiếm full-width với search real-time (client-side filtering).
  \item Mỗi job card hiển thị: tiêu đề, tên người đăng, khoa, địa điểm, mức lương,
    badge ``Khẩn cấp'' (nếu có), số ứng viên và thời gian đăng.
  \item Nút ``Load more'' theo kiểu infinite scroll giả -- fetch thêm 10 records từ Firestore.
\end{itemize}

\subsubsection{Màn hình chi tiết công việc (JobDetail)}

\begin{itemize}
  \item Header: tiêu đề task, badge danh mục, badge khẩn cấp.
  \item Thông tin người đăng: avatar, tên, khoa, điểm uy tín.
  \item Chi tiết: mô tả đầy đủ, kỹ năng yêu cầu (tags), địa điểm, thời gian, mức hỗ trợ.
  \item Nút ``Ứng tuyển ngay'' mở form nội tuyến (inline apply form) ngay trên trang -- không
    chuyển trang. Form có textarea cho thư giới thiệu và guard chống ứng tuyển trùng lặp.
\end{itemize}

\subsubsection{Màn hình hồ sơ cá nhân (Profile)}

\begin{itemize}
  \item Cột trái: avatar (Google OAuth), tên, khoa, MSSV, điểm uy tín và nút
    \textbf{``Xuất chứng nhận CV Passport''} màu xanh lá điều hướng đến \texttt{/cv-export}.
  \item Cột phải: tab lịch sử hoạt động (Tất cả / Đã hoàn thành / Đánh giá từ người thuê).
  \item Editable inline: skills (tags), bio.
\end{itemize}

\subsubsection{Màn hình xuất CV Passport (/cv-export)}

Giao diện chia 2 cột (tỷ lệ 3:2):
\begin{itemize}
  \item \textbf{Cột trái:} Preview chứng nhận live với thanh xanh trên/dưới, logo, tên sinh
    viên, 3 stat box (công việc / đánh giá / giờ làm), danh sách dự án.
  \item \textbf{Cột phải:} Panel ``Tùy chọn xuất file'' -- nút ``Tải xuống PDF'' (emerald),
    nút ``Chia sẻ LinkedIn'' (outline), toggle ``Hiển thị chi tiết đánh giá'', ghi chú xanh.
\end{itemize}

\subsubsection{Màn hình đăng nhập (Login)}

Trang login đơn giản với một nút duy nhất:
\textbf{``Đăng nhập bằng Google (Email Trường)''} -- thực hiện Google OAuth qua Firebase Auth,
gắn với domain \texttt{@hcmut.edu.vn} để xác thực danh tính.

\subsection{Responsive và Accessibility}

\begin{itemize}
  \item Tất cả layout dùng Tailwind CSS responsive breakpoints (\texttt{sm:}, \texttt{md:},
    \texttt{lg:}); trên mobile sidebar bộ lọc collapse thành bottom sheet.
  \item Color contrast ratio đạt WCAG AA cho text trên nền trắng và nền xanh.
  \item Tất cả form input có \texttt{aria-label} và \texttt{placeholder} rõ ràng.
  \item Toast notifications (\texttt{react-hot-toast}) thông báo kết quả mọi hành động.
\end{itemize}

\subsection{Công cụ và quy trình thiết kế}

\begin{itemize}
  \item \textbf{Figma:} thiết kế wireframe và mockup đầy đủ (file ID: \texttt{pNvShcrSUdVkdjTm8Pm4Bj}).
  \item \textbf{Tailwind CSS v4:} utility-first CSS, CSS variables làm design tokens.
  \item \textbf{Lucide React:} bộ icon nhất quán (Briefcase, Star, Zap, MapPin...).
  \item \textbf{Framer Motion:} micro-animation cho Hero section, card hover.
\end{itemize}
